\documentclass{article}
\usepackage{amsmath,amssymb,amsthm}
\usepackage{algorithm}
\usepackage{algpseudocode}
\usepackage{geometry}
\geometry{margin=1in}
\newtheorem{theorem}{Theorem}
\newtheorem{lemma}{Lemma}
\newtheorem{assumption}{Assumption}
\newcommand{\gap}{\mathcal{G}}
\newcommand{\diam}{\mathrm{diam}}
\begin{document}

\section*{Algorithm Name}
\textbf{Non‑Convex Frank–Wolfe (Conditional Gradient) Method}

\section*{Mathematical Formulation}
Let $f:\mathbb{R}^{d}\rightarrow\mathbb{R}$ be a continuously differentiable function and let 
\[
\mathcal{C}\subset\mathbb{R}^{d}
\]
be a non‑empty compact convex set (the feasible domain).  
Given a stepsize sequence $\{\gamma_{k}\}_{k\ge 0}\subset (0,1]$, the algorithm iterates as follows:

\begin{enumerate}
    \item \textbf{Linear minimization oracle (LMO)}  
    \[
    s_{k}\;:=\;\underset{s\in \mathcal{C}}{\arg\min}\;\big\langle \nabla f(x_{k}),\,s\big\rangle .
    \tag{1}
    \]
    \item \textbf{Convex combination update}  
    \[
    x_{k+1}\;:=\;(1-\gamma_{k})\,x_{k}\;+\;\gamma_{k}\,s_{k}.
    \tag{2}
    \]
\end{enumerate}
Define the \emph{Frank–Wolfe (FW) gap} at $x$ by
\[
\gap(x)\;:=\;\max_{s\in\mathcal{C}}\big\langle \nabla f(x),\,x-s\big\rangle
            \;=\; \big\langle \nabla f(x),\,x-s^{\star}(x)\big\rangle ,
\tag{3}
\]
where $s^{\star}(x)$ is any maximizer of the RHS. By construction of $s_{k}$ we have
\[
\big\langle \nabla f(x_{k}),\,s_{k}-x_{k}\big\rangle
   = -\gap(x_{k}).
\tag{4}
\]

\section*{Pseudocode}
\begin{algorithm}[H]
\caption{Non‑Convex Frank–Wolfe}
\begin{algorithmic}[1]
\Require Initial point $x_{0}\in\mathcal{C}$, stepsizes $\{\gamma_{k}\}_{k\ge0}$  
\For{$k=0,1,2,\dots$}
    \State Compute gradient $g_{k}\gets\nabla f(x_{k})$
    \State \textbf{LMO:} $s_{k}\gets\displaystyle\arg\min_{s\in\mathcal{C}}\langle g_{k},s\rangle$
    \State Update $x_{k+1}\gets (1-\gamma_{k})x_{k}+\gamma_{k}s_{k}$
    \If{termination criterion satisfied (e.g., $\gap(x_{k})\le\varepsilon$)}
        \State \textbf{break}
    \EndIf
\EndFor
\end{algorithmic}
\end{algorithm}

\section*{Dry Run Example}
Consider the one‑dimensional problem
\[
f(w) = (w-3)^{2},\qquad \mathcal{C}=[0,\,5],
\]
with $w_{0}=0$ and a constant stepsize $\gamma_{k}=0.1$.

\begin{center}
\begin{tabular}{c|c|c|c|c}
$k$ & $w_{k}$ & $\nabla f(w_{k})$ & $s_{k}$ (LMO) & $w_{k+1}$ \\ \hline
0 & $0$ & $2(0-3)=-6$ & $\displaystyle\arg\min_{s\in[0,5]}\{-6s\}=5$ & $0.9\cdot0+0.1\cdot5=0.5$\\[4pt]
1 & $0.5$ & $2(0.5-3)=-5$ & $\displaystyle\arg\min_{s\in[0,5]}\{-5s\}=5$ & $0.9\cdot0.5+0.1\cdot5=0.95$\\[4pt]
2 & $0.95$ & $2(0.95-3)=-4.1$ & $\displaystyle\arg\min_{s\in[0,5]}\{-4.1s\}=5$ & $0.9\cdot0.95+0.1\cdot5=1.355$\\[4pt]
3 & $1.355$ & $2(1.355-3)=-3.29$ & $\displaystyle\arg\min_{s\in[0,5]}\{-3.29s\}=5$ & $0.9\cdot1.355+0.1\cdot5=1.7195$
\end{tabular}
\end{center}
The objective values $f(w_{k})$ decrease from $9$ to $2.96$ after three iterations, illustrating the descent property of the method.

\newpage
\section*{Assumptions}
\begin{assumption}[Smoothness]\label{ass:smooth}
$f$ has $L$‑Lipschitz continuous gradient on $\mathcal{C}$, i.e.
\[
\|\nabla f(x)-\nabla f(y)\|\le L\|x-y\|,\qquad\forall x,y\in\mathcal{C}.
\tag{A1}
\]
\end{assumption}
\begin{assumption}[Compact Feasible Set]\label{ass:compact}
$\mathcal{C}$ is convex, compact and has finite diameter
\[
\diam(\mathcal{C})\;:=\;\max_{x,y\in\mathcal{C}}\|x-y\|\;<\;\infty .
\tag{A2}
\]
We denote $D:=\diam(\mathcal{C})$.
\end{assumption}
\begin{assumption}[Bounded Gradient on $\mathcal{C}$]\label{ass:gradbound}
There exists $G>0$ such that $\|\nabla f(x)\|\le G$ for all $x\in\mathcal{C}$.
\tag{A3}
\end{assumption}
No convexity of $f$ is required; the method is analyzed for possibly non‑convex $f$.

\section*{Main Convergence Theorem}
\begin{theorem}[Convergence of Non‑Convex Frank–Wolfe]\label{thm:main}
Under Assumptions~\ref{ass:smooth}–\ref{ass:gradbound}, let $\{\gamma_{k}\}_{k=0}^{K-1}$ be a constant stepsize
\[
\gamma_{k}\equiv \gamma\in(0,1].
\]
Define $f^{*}:=\inf_{x\in\mathcal{C}}f(x)$. Then for any $K\ge1$,
\[
\min_{0\le k\le K-1}\gap(x_{k})
\;\le\;
\frac{2\bigl(f(x_{0})-f^{*}\bigr)}{\gamma K}
\;+\;
\frac{L D^{2}}{2}\,\gamma .
\tag{5}
\]
Choosing the stepsize
\[
\gamma\;=\;\sqrt{\frac{2\bigl(f(x_{0})-f^{*}\bigr)}{L D^{2}K}}
\tag{6}
\]
optimizes the right‑hand side of (5) and yields the rate
\[
\min_{0\le k\le K-1}\gap(x_{k})
\;\le\;
\mathcal{O}\!\left(\frac{1}{\sqrt{K}}\right).
\tag{7}
\]
\end{theorem}

\section*{Theoretical Properties}
\begin{itemize}
    \item \textbf{Stability.} The recursion (5) depends only on $L$, $D$, and the initial suboptimality $f(x_{0})-f^{*}$; thus the method is robust to the choice of the initial point.
    \item \textbf{Stepsize Sensitivity.} Equation (5) shows a trade‑off: a larger $\gamma$ accelerates the linear term $2(f(x_{0})-f^{*})/(\gamma K)$ but inflates the quadratic term $LD^{2}\gamma/2$. The optimal $\gamma$ in (6) balances these effects.
    \item \textbf{Comparison to Gradient Descent.} For a smooth non‑convex problem without constraints, projected gradient descent attains a $O(1/\sqrt{K})$ bound on $\|\nabla f\|$. The FW gap $\gap(x_{k})$ plays an analogous role for constrained domains, and the rate (7) matches that of projected methods while avoiding costly projections.
\end{itemize}

\section*{Discussion}
The theorem guarantees that at least one iterate among the first $K$ steps has a small FW gap, i.e., it is an \emph{approximate stationary point} for the constrained problem. Because the algorithm only requires solving a linear problem over $\mathcal{C}$, it is particularly attractive when projections are expensive but a linear minimization oracle is cheap (e.g., the simplex, $\ell_{1}$-ball, or trace‑norm ball).

\newpage
\section*{Preliminaries (Definitions and Lemmas)}
\begin{lemma}[Descent Lemma]\label{lem:descent}
If $f$ satisfies Assumption~\ref{ass:smooth}, then for any $x,y\in\mathcal{C}$,
\[
f(y)\le f(x)+\langle\nabla f(x),y-x\rangle+\frac{L}{2}\|y-x\|^{2}.
\tag{8}
\]
\end{lemma}
\begin{proof}
Standard consequence of $L$‑Lipschitz gradient; see e.g. Nesterov (2004). \hfill $\square$
\end{proof}

\section*{Main Proof (Step‑by‑Step Derivation)}
We prove Theorem~\ref{thm:main}. Throughout the proof we keep the notation $g_{k}:=\nabla f(x_{k})$.

\bigskip
\noindent\textbf{[Step 1]} Apply Lemma~\ref{lem:descent} with $x=x_{k}$ and $y=x_{k+1}$:
\[
f(x_{k+1})\le f(x_{k})+\langle g_{k},x_{k+1}-x_{k}\rangle+\frac{L}{2}\|x_{k+1}-x_{k}\|^{2}.
\tag{9}
\]

\noindent\textbf{[Step 2]} Substitute the update rule (2):
\[
x_{k+1}-x_{k}= \gamma_{k}(s_{k}-x_{k}),
\qquad
\|x_{k+1}-x_{k}\|^{2}= \gamma_{k}^{2}\|s_{k}-x_{k}\|^{2}.
\tag{10}
\]

\noindent\textbf{[Step 3]} Insert (10) into (9):
\[
f(x_{k+1})\le f(x_{k})+\gamma_{k}\langle g_{k},s_{k}-x_{k}\rangle
        +\frac{L}{2}\gamma_{k}^{2}\|s_{k}-x_{k}\|^{2}.
\tag{11}
\]

\noindent\textbf{[Step 4]} By definition of the LMO (1) we have (4):
\[
\langle g_{k},s_{k}-x_{k}\rangle = -\gap(x_{k}).
\tag{12}
\]

\noindent\textbf{[Step 5]} Use the diameter bound $ \|s_{k}-x_{k}\|\le D $ (from Assumption~\ref{ass:compact}):
\[
\|s_{k}-x_{k}\|^{2}\le D^{2}.
\tag{13}
\]

\noindent\textbf{[Step 6]} Combine (11), (12) and (13):
\[
f(x_{k+1})\le f(x_{k}) - \gamma_{k}\,\gap(x_{k}) + \frac{L D^{2}}{2}\,\gamma_{k}^{2}.
\tag{14}
\]

\noindent\textbf{[Step 7]} Rearrange (14) to isolate the gap term:
\[
\gamma_{k}\,\gap(x_{k})
\le f(x_{k})-f(x_{k+1}) + \frac{L D^{2}}{2}\,\gamma_{k}^{2}.
\tag{15}
\]

\noindent\textbf{[Step 8]} Sum inequality (15) over $k=0,\dots,K-1$:
\[
\sum_{k=0}^{K-1}\gamma_{k}\,\gap(x_{k})
\le f(x_{0})-f(x_{K}) + \frac{L D^{2}}{2}\sum_{k=0}^{K-1}\gamma_{k}^{2}.
\tag{16}
\]

\noindent\textbf{[Step 9]} Since $f(x_{K})\ge f^{*}$, replace $-f(x_{K})$ by $-f^{*}$:
\[
\sum_{k=0}^{K-1}\gamma_{k}\,\gap(x_{k})
\le f(x_{0})-f^{*} + \frac{L D^{2}}{2}\sum_{k=0}^{K-1}\gamma_{k}^{2}.
\tag{17}
\]

\noindent\textbf{[Step 10]} For a constant stepsize $\gamma_{k}\equiv\gamma$, the left‑hand side becomes $\gamma\sum_{k=0}^{K-1}\gap(x_{k})$ and the right‑hand side simplifies:
\[
\gamma\sum_{k=0}^{K-1}\gap(x_{k})
\le f(x_{0})-f^{*} + \frac{L D^{2}}{2}\,K\gamma^{2}.
\tag{18}
\]

\noindent\textbf{[Step 11]} Divide both sides of (18) by $\gamma K$:
\[
\frac{1}{K}\sum_{k=0}^{K-1}\gap(x_{k})
\le \frac{f(x_{0})-f^{*}}{\gamma K}
   + \frac{L D^{2}}{2}\,\gamma .
\tag{19}
\]

\noindent\textbf{[Step 12]} Since the minimum of a set is bounded by its average,
\[
\min_{0\le k\le K-1}\gap(x_{k})
\le \frac{1}{K}\sum_{k=0}^{K-1}\gap(x_{k}),
\tag{20}
\]
and combining (19)–(20) yields exactly inequality (5) of the theorem.

\noindent\textbf{[Step 13]} To obtain the optimal constant stepsize, minimize the RHS of (5) w.r.t.\ $\gamma$:
\[
\phi(\gamma)=\frac{2\Delta}{\gamma K}+\frac{L D^{2}}{2}\gamma,
\qquad\Delta:=f(x_{0})-f^{*}.
\]
Setting $\phi'(\gamma)=0$ gives
\[
-\frac{2\Delta}{\gamma^{2}K}+\frac{L D^{2}}{2}=0
\;\Longrightarrow\;
\gamma^{2}= \frac{4\Delta}{L D^{2}K},
\]
hence the optimal choice (6). Substituting (6) into (5) yields the $O(1/\sqrt{K})$ bound (7). \hfill $\square$

\section*{Rate Derivation (With Constants)}
From (6) we have
\[
\gamma = \sqrt{\frac{2\Delta}{L D^{2}K}} .
\]
Plugging into (5):
\[
\min_{k}\gap(x_{k})
\le \frac{2\Delta}{\gamma K} + \frac{L D^{2}}{2}\gamma
= \frac{2\Delta}{K}\sqrt{\frac{L D^{2}K}{2\Delta}} + \frac{L D^{2}}{2}\sqrt{\frac{2\Delta}{L D^{2}K}}
= 2\sqrt{\frac{L D^{2}\Delta}{2K}} + \sqrt{\frac{L D^{2}\Delta}{2K}}
= 3\sqrt{\frac{L D^{2}\Delta}{2K}} .
\]
Thus,
\[
\boxed{\displaystyle
\min_{0\le k\le K-1}\gap(x_{k})\;\le\;3\sqrt{\frac{L D^{2}\bigl(f(x_{0})-f^{*}\bigr)}{2K}}
= \mathcal{O}\!\left(\frac{1}{\sqrt{K}}\right)} .
\]

\section*{Special Cases}
\begin{itemize}
    \item \textbf{Convex $f$.} When $f$ is convex, the FW gap upper bounds the duality gap, and the same $O(1/K)$ rate can be recovered by using the diminishing stepsize $\gamma_{k}=2/(k+2)$.
    \item \textbf{Polytope $\mathcal{C}$.} If $\mathcal{C}$ is a polytope, the LMO reduces to a linear program that returns an extreme point; the bound (5) still holds unchanged.
    \item \textbf{Exact line‑search.} Replacing the constant $\gamma$ by the exact line‑search
    $\gamma_{k}^{\star}= \arg\min_{\gamma\in[0,1]} f\bigl((1-\gamma)x_{k}+\gamma s_{k}\bigr)$
    can only improve the RHS of (14), thus the $O(1/\sqrt{K})$ guarantee still applies.
\end{itemize}

\end{document}